\begin{document}

{
\setcounter{tocdepth}{2}
\tableofcontents
}
\section{TFPT / Universalraum: Zusammensetzung aus der gemeinsamen
Antwort}\label{tfpt-universalraum-zusammensetzung-aus-der-gemeinsamen-antwort}

Vollständige Fortschreibung v1.6.9 · 15. September 2026 · Bezugsfassung
v1.6.8

\subsection{1. Ergebnis: Die fehlende Verbindung wird berechenbar, aber
noch nicht ursprünglich
ausgewählt}\label{ergebnis-die-fehlende-verbindung-wird-berechenbar-aber-noch-nicht-urspruxfcnglich-ausgewuxe4hlt}

Die drei beauftragten Arbeitslinien wurden parallel ausgeführt:
ursprüngliche Komposition, ein kompletter kleiner Ablauf und ein früher
relativistischer Feldtest. Hinzu kommt eine unabhängige globale
CAR-Rechnung im Hauptstrang. Sie ergeben keinen Abschluss der TOE, aber
vier konkrete Fortschritte:

\begin{enumerate}
\def\labelenumi{\arabic{enumi}.}
\tightlist
\item
  \textbf{Gemeinsame Ressourcen sind im ursprünglichen Modell bereits
  vorhanden.} Es gibt 64 gemeinsam benutzte Fermionmoden, nicht 480
  unabhängige Bankbausteine. Jeder Fermionmodus tritt in 15 Vertices
  auf. Was noch fehlt, ist die Quellregel für mehrere vollständige,
  miteinander überlappende Vorkommen dieses Trägers.
\item
  \textbf{Die dynamische Antwort kann die zusätzliche Zusammensetzung
  zurückliefern.} Ein vierter Energiemoment bestimmt in einer kleinen
  Zwei-Chart-Familie ihre Überlappung. In einer anderen normierten
  Familie rekonstruiert die matrixaufgelöste Dreiladungsantwort die
  positive Kopplungsmatrix. Letztere unterscheidet sogar Modelle, deren
  gesamte Zweiladungsspektren gleich sind.
\item
  \textbf{Der Eingriff wurde einfacher und die Aufzeichnung vollständig
  bilanziert.} Der bisher zusätzlich gewährte Einzelmodenimpuls wird
  durch die Bosonzahlphase ersetzt. Ein kohärenter zweistufiger Zeiger
  liefert einen wirklichen aufgezeichneten Prozess, erhält alle
  Paarkohärenzen und erzeugt einen exakt berechneten unbedingten
  Folgeeffekt. Kontrolle, Zeiger und Arbeit bleiben ausdrücklich
  zusätzliche Ressourcen.
\item
  \textbf{Der relativistische Ausleseadapter wurde konstruiert und früh
  gegen Kinetik getestet.} Ein Impulsfaktor genügt für den
  gemischthändigen algebraischen Kanal. Der freie masselose
  Produktansatz liefert damit allerdings exakt null. Die Feldzuordnung
  muss außerdem den ursprünglichen Vertex einschließlich Adjunktion
  erhalten.
\end{enumerate}

Die leitende Vereinfachung lautet jetzt: \textbf{Nicht möglichst viele
ähnliche Zahlen sammeln, sondern die Zusammensetzungsregel an einer
höheren Antwort eindeutig unterscheidbar machen.} Das ersetzt keine
Quellherleitung. Es liefert ihr erstmals in diesen Familien einen
expliziten Rückrechnungs- und Falsifikationstest.

„Exakt`` bezeichnet in dieser Ausgabe algebraische Identitäten und
vollständig endliche Rechnungen im angegebenen Vertrag. „Bedingt``
bezeichnet zusätzlich vorausgesetzte Träger, Präparationen, Kontrollen
oder Feldzuordnungen. Die Programme sind Reproduktionshilfen, keine
Lean-Formalisierung. Die vollständige v1.6.8 bleibt unverändert als
historischer Teil enthalten; maßgeblich für den aktuellen Status sind
diese Kapitel und die neuen vollständigen Teilberichte.

\subsection{2. Einfaches Gesamtbild}\label{einfaches-gesamtbild}

Man kann sich den ursprünglichen Tensor wie eine \textbf{Kupplung}
vorstellen: Zwei Fermionen können gemeinsam in ein Boson übergehen und
zurück. Diese Kupplung ist intern wesentlich besser bestimmt, als es die
offenen Grundsatzfragen vermuten lassen. Die offene Frage lautet aber
noch: Welche Kupplungen greifen auf \textbf{dieselben} Ressourcen zu,
und welche sind unabhängig?

Zwei Maschinen können beim einfachen Anschlagen gleich klingen und beim
nächsten Versuch verschieden reagieren. Genau das zeigt die neue
N=2-/N=3-Rechnung. Das höhere Antwortmuster legt innerhalb der erklärten
Familie die zusätzliche Kopplungsmatrix offen. Die Metapher der Schatten
wird hier damit zu einer konkreten inversen Aufgabe.

{\def\LTcaptype{none} % do not increment counter
\begin{longtable}[]{@{}
  >{\raggedright\arraybackslash}p{(\linewidth - 4\tabcolsep) * \real{0.3333}}
  >{\raggedright\arraybackslash}p{(\linewidth - 4\tabcolsep) * \real{0.3333}}
  >{\raggedright\arraybackslash}p{(\linewidth - 4\tabcolsep) * \real{0.3333}}@{}}
\toprule\noalign{}
\begin{minipage}[b]{\linewidth}\raggedright
Gesuchtes gemeinsames Bild
\end{minipage} & \begin{minipage}[b]{\linewidth}\raggedright
Jetzt tatsächlich berechnet
\end{minipage} & \begin{minipage}[b]{\linewidth}\raggedright
Noch fehlende Brücke
\end{minipage} \\
\midrule\noalign{}
\endhead
\bottomrule\noalign{}
\endlastfoot
Eine ursprüngliche Regel & W, gemeinsamer CAR-Träger, Clock und innere
Symmetrie & Auswahl globaler Einbettungen und ihres Hamiltonoperators \\
Verschiedene Schatten desselben Ablaufs & Exakte höhere Antworten und
zwei inverse Formeln & Compilerseitig zugängliche Präparation und
Messung dieser Antworten \\
Information, Eingriff und Zeit & Ein gemeinsamer kleiner kontrollierter
Prozess samt Zeiger & Autonome Quelle von Kontrollen, Uhr, Zeiger und
Steuerarbeit \\
Physikalische Teilchen & Korrekte endliche Ladungen und bedingte
Lorentz-Auslese & Nichtverschwindende Kinetik, native Massenschale und
chirales Maß \\
Unsere Raumzeit & Klare Tests für einen künftigen wachsenden Träger &
3+1D-Lokalität, relativistischer Limes und dynamische Gravitation \\
\end{longtable}
}

Das ist noch keine holographische Dualität. Eine solche stärkere
Behauptung müsste insbesondere festlegen, welche Rand- und inneren
Observablen einander entsprechen, wie ihre gemeinsame Dynamik erhalten
bleibt und welche Daten vollständig sind. Die hier bewiesene
Teiltraceformel ist eine endliche inverse Identität unter klaren
Annahmen, kein Flächengesetz und kein Beweis, dass Realität ein Rand
dieses Modells ist.

\subsection{3. Quellenweg: ursprünglicher Clock ist keine
Kopiermaschine}\label{quellenweg-urspruxfcnglicher-clock-ist-keine-kopiermaschine}

Die alten Quellen wurden bis zu ihren tatsächlichen Trägern
zurückverfolgt. Die sieben frühen Wörter wirken auf demselben C3; ihre
49 Produkte schließen dort eine siebendimensionale Algebra. Die 60
ursprünglichen Gaußwurzelstrahlen wirken auf demselben C4. Der spätere
W-Konstruktor liefert einen gemeinsamen 64-Fermionträger und 60
Bosonkanäle. Keiner dieser Schritte darf allein durch neue Indexnamen zu
Raumorten oder unabhängigen Fockkopien werden.

Für den dokumentierten passiven Clock gilt exakt mit dem fermionischen
Paarlift L und dem bosonischen Lift G\_B:

\[\adjustbox{max width=\linewidth}{$\displaystyle 
L^kW^T=W^TG_B^k,\qquad k=0,\ldots,5.
$}\]

Alle sechs Clockbilder des hellen Raums sind daher derselbe Raum. Ihr
Vereinigungsrang ist \textbf{60, nicht 360}. Diese Aussage wurde durch
exakte Matrixidentitäten kontrolliert, nicht durch numerische Ränge. Ein
anderer aktiver Clockvertrag wäre gesondert zu konstruieren.

Der Quellenstrang untersucht dann zwei isometrisch eingebettete
Fermion-Charts

\[\adjustbox{max width=\linewidth}{$\displaystyle 
f_i^L=a_i,\qquad f_i^R=c\,a_i+\sqrt{1-c^2}\,d_i.
$}\]

Jeder trägt exakt den ursprünglichen W mit derselben lokalen Stärke g
und besitzt eine eigene Bosonbank. Gegenseitig gilt
\(\{f_i^L,f_j^{R\dagger}\}=c\delta_{ij}\). Die Charts sind überlappende
Operationsbereiche, \textbf{keine schon hergeleiteten Raumorte}.

Der vollständige Ablauf von einem linken Boson zur rechten Bosonzahl
schließt für jeden inneren Kanal in fünf Zuständen; vom gegebenen
Eingang sind vier zyklisch erreichbar. Die Anfangswahrscheinlichkeit ist

\[\adjustbox{max width=\linewidth}{$\displaystyle 
P(L\to R;t)=16g^4c^4t^4+O(t^6).
$}\]

Mit gleichmäßig über alle 60 linken Bosonen gemischtem Eingang und
Gesamt-Bosonzahl rechts ist dieser Vergleich voll G-invariant. Die zwei
rationalen Überlappungen c=3/5 und c=4/5 erhalten sämtliche einzelnen
lokalen W-Identitäten, liefern aber verschiedene Übergänge. Die
Präparations- und Detektorverfügbarkeit wird nicht aus G-Invarianz
gefolgert.

Für \(\mu_k=\operatorname{tr}(\rho_LH^k)\) ergibt sich die skalenfreie
inverse Formel

\[\adjustbox{max width=\linewidth}{$\displaystyle 
\boxed{I_4=\frac{\mu_4-3\mu_1^2\mu_2+2\mu_1^4}{(\mu_2-\mu_1^2)^2}=1+c^4.}
$}\]

Die ersten drei Momente sind gleich; das vierte unterscheidet die
Modelle. Für diese nichtnegative Familie ist c aus I4 eindeutig. Keine
Vermutung über Primzahlen oder eine passende Dezimalzahl wird dafür
benötigt.

Die ursprünglichen C4-Strahlen enthalten tatsächlich
Überlappungsquadrate \(0,1/4,1/2,1\). Das ist ein sinnvoller nächster
Ansatzpunkt. Es fehlt jedoch ihre \textbf{typisierte Zuordnung zu
globalen CAR-Einbettungen}. Der bereits interne SU(4)-Fundamentalfaktor
ist nicht automatisch ein neuer G-trivialer Kopienfaktor. Weder die
bekannten Überlappungszahlen noch die innere Symmetrie allein wählen c,
die Bosonbanken und die globale Vertexregel aus.

\subsection{4. Stärkerer Vergleich: selbst gleiche gesamte
Zweiladungsspektren reichen
nicht}\label{stuxe4rkerer-vergleich-selbst-gleiche-gesamte-zweiladungsspektren-reichen-nicht}

Eine unabhängige Familie verwendet nur eine gemeinsame Bosonbank und
einen zusätzlichen symmetrischen normierten Tensor C auf L
Fermionkopien. Ihre globalen Paarmatrizen sind \(M_A=C\otimes A_A\). In
dieser Familie bleiben das globale Paargramm \(8I_{60}\) und das gesamte
N=2-Hamiltonspektrum unverändert. Die lokalen Einzelkopienamplituden
bleiben dagegen nicht unverändert; diese Annahme ist ausdrücklich
verschieden von Abschnitt 3.

Die verteilte Komposition \(C=I_L/\sqrt L\) und die kollektive
Komposition \(C=J_L/L\) erfüllen dieselben Norm- und
Permutationsbedingungen. Schon bei L=2 besitzen sie dennoch
unterschiedliche gesamte N=3-Spektren. Bei Energie Δ gibt es im
verteilten Modell keinen Eigenzustand, im kollektiven genau 64, für
Δ\textgreater0 und g≠0.

Das ist ein Unterschied der vollständigen Energieeigenwerte mit
Multiplizitäten, nicht bloß verschiedener Portnamen. Er wurde durch
vollständige CAR-Kontraktion auf dem 349056-dimensionalen N=3-Träger und
exakt kontrollierte kleinere Gramoperatoren geprüft.

Die allgemeine Formel und ihre Umkehrung lauten, mit \(Q=C^\dagger C\)
und dem nur aus W gebauten K:

\[\adjustbox{max width=\linewidth}{$\displaystyle 
G_C^{(3)}=8I-Q\otimes K,\qquad\operatorname{tr}K=960,
$}\] \[\adjustbox{max width=\linewidth}{$\displaystyle 
\boxed{Q=\frac1{960}\operatorname{Tr}_{\rm intern}(8I-G_C^{(3)}).}
$}\]

Hier wird der \textbf{matrixaufgelöste} Gramoperator einschließlich
Tensorfaktor- und Phasendictionary benötigt. Eigenwerte allein geben
nicht die Orientierung von Q in einer markierten Kopienbasis. Die
Gleichung ist kein Beweis, dass der Compiler Q auswählt. Sie ist ein
exakter Test, welche Zusammensetzung eine gegebene höhere Antwort trägt.

Der Anschluss zur konkreten Bosonzahlmessung ist ebenfalls berechnet:
Für einen bosonseitigen normierten Eingang v folgt

\[\adjustbox{max width=\linewidth}{$\displaystyle 
p_{N_b=1}(t)=1-g^2\langle v,G_C^{(3)}v\rangle t^2+O(t^4).
$}\]

Mit hinreichend zugänglichen Eingängen könnte dieselbe Art von Auslese
die Matrix rekonstruieren. Die benötigten Eingänge, Zeitkalibrierung und
Messungen sind zusätzliche Bedingungen, nicht schon vorhandene
Compilergeräte.

Ein einzelner symmetrischer C-Tensor besitzt zudem immer eine
Takagi-Basis mit erhaltenen einzelnen Fermionkopienparitäten. Das ist
eine genaue Grenze dieses besonders einfachen Ansatzes, keine allgemeine
Grenze von TFPT. Zwei nichtkompatible Chart-Kopplungen wie in Abschnitt
3 haben diese spezielle gemeinsame Zerlegung nicht. Damit ist auch klar,
welche zusätzliche Struktur für einen echten Einfermiontransfer
mindestens geprüft werden sollte.

\subsection{5. Vereinfachter vollständiger Prozess mit einer echten
Aufzeichnung}\label{vereinfachter-vollstuxe4ndiger-prozess-mit-einer-echten-aufzeichnung}

Der native Dreizustandsraum
\(\mathcal K_3=\operatorname{span}(p_0,R_7,b_0)\) aus v1.6.8 bleibt
unverändert. Neu ist die mittlere Operation:

\[\adjustbox{max width=\linewidth}{$\displaystyle 
Z_b=e^{i\pi N_b}.
$}\]

Auf K3 liegt sie in der bereits dokumentierten Kontrollalgebra aus X und
Nb, komplexe Dimension fünf. Der alte Einzelmodenimpuls Z4 liegt nicht
darin. Ein separat benötigter Modenselektor entfällt also tatsächlich.

Nach einer Vorentwicklung, dem optionalen Bosonimpuls und derselben
Nachentwicklung ändert sich die unbedingte Besetzung der Fermionmode 5.
Am Prüfpunkt g/Δ=1/20:

{\def\LTcaptype{none} % do not increment counter
\begin{longtable}[]{@{}lr@{}}
\toprule\noalign{}
Protokoll & Unbedingte Wahrscheinlichkeit n5=1 \\
\midrule\noalign{}
\endhead
\bottomrule\noalign{}
\endlastfoot
Ohne mittleren Eingriff & 0.0002194998817122665 \\
Mit Bosonzahlphase & 0.0005299169088385700 \\
Differenz & +0.0003104170271263035 \\
\end{longtable}
}

Der Sendergriff und der Empfängeroperator kommutieren zum
Eingriffszeitpunkt; erst die weitere native Entwicklung vermittelt den
Unterschied. Das ist ein Eingriffstest, keine bloße Korrelation. Anders
als im alten Versuch endet der Impulsarm nicht bosonfrei. Auch ist Nb
ein globaler Bosongriff, kein bereits räumlich lokaler Sender.

Ein echter Recorder koppelt einen bereiten neutralen Zweizustandszeiger
nur an die Unterscheidung Paar/Boson:

\[\adjustbox{max width=\linewidth}{$\displaystyle 
V\psi=(I-N_b)\psi\otimes|0\rangle+N_b\psi\otimes|1\rangle.
$}\]

Diese Form gilt auf dem verwendeten Sektor N=2. Der gemeinsame Träger
hat sechs Dimensionen. Alle Paare behalten dasselbe Zeigerlabel, ihre
interne Interferenz bleibt erhalten. Wird der Zeiger nicht ausgewählt,
sondern ignoriert, ergibt sich der echte CPTP-Update
\((\rho+Z_b\rho Z_b)/2\) und genau der halbe unbedingte Folgeeffekt.
Feine Paarlabels zu messen und anschließend umzubenennen wäre ein
anderer, kohärenzzerstörender Prozess.

Die Ressourcen werden nicht mehr versteckt: Der Phasenimpuls deponiert
im System die Energie \textbf{Δ/54}, der Recorder \textbf{Δ/108}. Beide
erhalten die Gesamtladung. Der Controller und der Zeiger müssen jedoch
vorbereitet, gekoppelt und energetisch berücksichtigt werden. Aus einem
festen H folgt keine unabhängige X-/Nb-Schaltbarkeit; aus einer
verfügbaren Phase folgt kein kontrollierter Zeigerkoppler.

Eine zusätzliche allgemeine Grenze lautet: Wenn ein echter innerer
Symmetrieimpuls V sowohl mit H als auch dem anfänglichen Empfänger O
kommutiert, kommutiert er mit O(t) zu jeder späteren Zeit. Er kann
diesen verzögerten Eingriffseffekt deshalb nicht ersetzen.
\textbf{G-invariant} ist etwas anderes als \textbf{Element der inneren
Gruppe G}. Zb erfüllt das Erste, nicht das Zweite.

\subsection{6. Relativistischer Anschluss: ein positiver Adapter und ein
negativer
Kinetiktest}\label{relativistischer-anschluss-ein-positiver-adapter-und-ein-negativer-kinetiktest}

Für ein gemischthändiges Komposit \(T\in(1,1/2)\) ist ohne zusätzliche
Lorentzstruktur keine lineare Auslese in einen linken Weylspinor
möglich. Mit genau einem Impulsfaktor P gibt es dagegen einen bis auf
Normierung eindeutigen solchen Operator C\_P. Für jedes P≠0 hat er Rang
zwei und einen vierdimensionalen Kern.

Eine explizite Einsetzung J\_P erfüllt

\[\adjustbox{max width=\linewidth}{$\displaystyle 
C_PJ_P=3\det(P)I_2.
$}\]

Bei zeitartigem P liefert das einen kovarianten Rückweg und Projektor.
Bei nullartigem P bleibt die Auslese bestehen, aber es gibt keinen nur
von P abhängigen kovarianten Rückweg auf den ganzen Weylraum. Ein
zusätzliches zeitartiges Referenz-Q ermöglicht einen Rückweg; seine
Herkunft wäre eine neue Voraussetzung. Diese Q-Referenz ist
\textbf{nicht} die Kopplungsmatrix Q aus Abschnitt 4.

Der entscheidende frühe Physiktest wurde nicht vertagt: Ist das Komposit
ein Produkt eines freien masselosen selbstdualen Feldes F und eines
freien masselosen Weylfeldes, dann gilt auf deren Bewegungsgleichungen
exakt

\[\adjustbox{max width=\linewidth}{$\displaystyle 
C_{P_F+P_\psi}(F\otimes\bar\psi)=0.
$}\]

Die Identität gilt auch für nichtkollineare Impulse und zeitartigen
Gesamtimpuls. Der algebraisch passende Adapter ergibt in diesem
konkreten freien Ansatz also \textbf{kein sichtbares Teilchen}. Ein
nichtnuller Kanal müsste aus einer wirklich berechneten Wechselwirkung,
Masse, anderen Kinetik oder anderen Feldzuordnung kommen.

Zusätzlich wurde der ursprüngliche Vertex geprüft: Wird f wörtlich als
linkes Weylfeld und b†ff als ableitungsfreier Lorentzskalar
interpretiert, muss das Feld b† links, sein Adjungiertes b rechts sein.
Dann ist D=bf† gleichhändig rechts, während χ†=b†ηf† gemischthändig ist.
Die Zuordnung vertauscht sich gegenüber der isolierten Annahme „b
links``. Nicht beide Konventionen dürfen gleichzeitig für denselben
Vertex benutzt werden.

Das ist kein Ausschluss einer konsistenten Moden-Feld-Abbildung: Ein
nativer Fock-Erzeuger ist noch kein vollständiges lokales Feld. Genau
diese Abbildung mit Wellenfunktionen, Ladungen und Adjungierten ist zu
konstruieren. Die Ladungen −1 und +3 bleiben verschieden; ein neutraler
Impulsfaktor setzt sie nicht gleich. Die Spinorkonventionen entsprechen
\href{https://arxiv.org/abs/0812.1594}{Dreiner, Haber und Martin}; die
konkreten hier genannten Matrix- und Nullidentitäten wurden selbst
geprüft.

